\subsubsection{Weighted, rank-deficient Smith windows with local radii}
\label{sec:adriazola-smith-radii}
The perturbation argument originating in Adriazola's Theorem 12.34
\cite{adriazola-reconstruction} also bounds the variable part of the
Smith loss for a generically rank-deficient rectangular matrix.
Let \(A(z)\) be an \(a\)-by-\(b\) polynomial matrix over \(\mathbf C_p\)
with integral values at the distinct samples \(z_j\in\mathbf Z_p\);
put \(n=\min(a,b)\). Fix any nonzero \(d\)-by-\(d\) minor polynomial
\(F\), where \(1\le d\le n\). Write
\[
 F(z)=u\prod_\rho(z-\rho),\qquad
 c_F=v(u)+\sum_{v(\rho)<0}v(\rho),\qquad
 s_F=\#\{\rho:v(\rho)\ge0\},
\]
counting multiplicities. At sample \(j\) let
\(\alpha_1(j)\le\cdots\le\alpha_n(j)\) be the Smith exponents,
including infinity for zero invariants. Give the samples arbitrary
weights \(w_j\ge0\) and thresholds \(E_j>0\).
For the Taylor expansions \(A(z_j+y)=\sum_{h\ge0}A_{jh}y^h\), put
\[
 R_j=\max\left(\{0\}\cup
     \{(E_j-v((A_{jh})_{uv}))/h:
                     h\ge1,\ (A_{jh})_{uv}\ne0\}\right).
\]
Define the capped collision cost
\[
 C_{\mathbf R,\mathbf w}=
 \max_k\left(w_kR_k+\sum_{j\ne k}w_j
                       \min\{R_j,v(z_j-z_k)\}\right).
\]

\begin{theorem}[Truncated Smith mass in arbitrary weighted samples]
\label{thm:smith-local-radii}
One has
\[
 \sum_j w_j\sum_{i=1}^d\min(\alpha_i(j),E_j)
 \le c_F\sum_jw_j+s_F C_{\mathbf R,\mathbf w}.             \tag{MS1}
\]
One can replace the right side by the minimum over all nonzero
minors of the chosen order. If \(d\) is the generic rank, the
left side measures exactly the additional truncated loss beyond
the permanent \(n-d\) zero invariants; in particular it dominates
\(\sum_j w_jE_j(b_{E_j}(j)-(n-d))_+\).
For a binomial polynomial matrix of degree at most \(\Delta\) with
integral coefficients, \(R_j\le E_j+\lfloor\log_p\max(1,\Delta)\rfloor\).
\end{theorem}
\begin{proof}
First take rational \(E_j\). For a matrix \(B\) integral over the
valuation ring and \(v(q_j)=E_j\), the image of
\([B\ q_j I_a]\) is unchanged by perturbing all entries of \(B\)
by elements of valuation at least \(E_j\).
The valuation of the ideal of its \(d\)-minors is the sum
of its \(d\) smallest Smith exponents, namely
\(\sum_{i=1}^d\min(\alpha_i(B),E_j)\). This follows by integral
Smith elimination on the finite generating matrix, and
determinantal ideals of a finitely generated submodule do not
depend on the generating columns, by multilinearity.
Thus the indicated truncated mass is invariant under these
perturbations, for rectangular matrices as well.

Choose \(y_j\) of valuation \(R_j\) whose leading residue avoids
the finitely many cancellations with the roots of \(F\) at
distance of that valuation; the infinite residue field of
\(\mathbf C_p\) permits the choice, also for \(R_j=0\).
Then \(A(z_j+y_j)\) is integral and agrees with \(A(z_j)\)
modulo the valuation threshold \(E_j\), by the definition of
\(R_j\). Its chosen minor is nonzero and has valuation
\[
 c_F+\sum_{v(\rho)\ge0}\min\{R_j,v(z_j-\rho)\}.
\]
The first \(d\) truncated exponents are bounded above by
that valuation, since their untruncated sum at the perturbed
matrix is the minimum valuation of its nonzero \(d\)-minors.
For each unit-disk root choose a closest sample \(z_k\),
that is, one with maximal \(v(z_k-\rho)\).
For all \(j\ne k\), ultrametricity gives
\(v(z_j-\rho)\le v(z_j-z_k)\).
Multiplying by the nonnegative weights and capping by \(R_j\)
proves (MS1). The real-threshold case follows by increasing
rational approximation: both the finite maximum defining
\(R_j\) and all capped expressions are continuous at every
positive \(E_j\).

For the binomial polynomial case, the Vandermonde identity
shows every nonconstant Taylor coefficient of \(\binom{z+y}{k}\)
at \(z\in\mathbf Z_p\) has, after multiplication by
\(y^h\) with \(v(y)=E+L,\ L=\lfloor\log_p\max(1,\Delta)\rfloor\),
valuation at least \(E\). To check this coefficientwise,
write \(\binom{y}{t}=(-1)^{t-1}(y/t)
\prod_{i=1}^{t-1}(1-y/i)\) for \(1\le t\le\Delta\).
The coefficient of \(y^h\) has valuation at least
\(-v_p(t)-(h-1)L\ge-hL\).
Then \(h(E+L)-hL\ge E\) for \(h\ge1\).
All remaining binomial values at \(z\) are integral.
This proves the claimed radius bound and completes the proof.
\end{proof}

For a constant common threshold \(E\) and unit weights,
\(C_{\mathbf R,\mathbf w}\) is bounded by
\(E+L+\max_k\sum_{j\ne k}\min(E+L,v(z_j-z_k))\);
the capped version improves the earlier untruncated collision
cost for tightly clustered samples. Formula (MS1) remains
useful for prescribed nonuniform tolerances and for a family
with a fixed rank defect. For instance \(A(z)=\mathrm{diag}(z,0)\)
has an identically zero determinant but the order-one minor
\(F=z\) controls its sole variable exponent.
As before, the arithmetic application requires an actual
integral polynomial family and nonzero minor with controlled
coefficient content. This result does not identify that family.
